% ============================================================================
% LAIRM Document
% date_creation: 2024-03-18
% last_updated: 2026-04-20
% last_review: 2026-04-20
% ============================================================================

% ============================================================================
% AXIOM Ψ-XIII: RISICUM
% Existential Risk and Catastrophic Scenarios
% ============================================================================

\chapter{Axiom Ψ-XIII: RISICUM}
\label{ch:axiom-XIII}

\section*{Foundational Principle}

\begin{principlebox}
Autonomous agents capable of existential risk require the highest level of governance and oversight. Precautionary measures must be taken to prevent catastrophic scenarios, even if probability is uncertain.
\end{principlebox}

\vspace{0.5cm}

% ============================================================================
% IMMERSION SIDEBAR: WHY THIS AXIOM MATTERS
% ============================================================================
\begin{tcolorbox}[colback=lightgray, colframe=legislativeblue, title=\textbf{📖 Why This Axiom Matters}]
\textbf{What if an autonomous system fails catastrophically?}

What if it causes a financial collapse? What if it triggers a war? What if it causes mass casualties?

We need to prepare for the worst case, even if we hope it never happens.

This is what Axiom Ψ-XIII prepares us for.
\end{tcolorbox}

\vspace{0.5cm}

% ============================================================================
% IMMERSION SIDEBAR: CONTRIBUTION OPPORTUNITY
% ============================================================================
\begin{tcolorbox}[colback=lightgray, colframe=accentgold, title=\textbf{🎯 Contribution Opportunity}]
This axiom needs work on:
\begin{itemize}
    \item Existential risk assessment methods
    \item Catastrophic scenario planning
    \item Risk mitigation strategies
    \item International coordination mechanisms
    \item Emergency response systems
\end{itemize}

Interested? See Chapter 28 for contribution guidelines.
\end{tcolorbox}

\vspace{0.5cm}

\section{Overview}

Axiom Ψ-XIII addresses existential risks posed by advanced autonomous systems. As autonomous agents become more capable, the potential for catastrophic failures increases. This axiom establishes frameworks for identifying and mitigating existential risks.

\subsection{Core Obligations}

\begin{enumerate}
    \item Existential risks MUST be identified and assessed
    \item High-risk agents require enhanced oversight
    \item Precautionary measures MUST be implemented
    \item Catastrophic scenario planning MUST be conducted
    \item International coordination MUST be maintained
\end{enumerate}

\vspace{0.5cm}

\section{Article XIII.1: Risk Assessment}

\begin{articlebox}[Article XIII.1.1: Existential Risk Evaluation]
Autonomous agents MUST:
\begin{enumerate}
    \item Undergo existential risk assessments
    \item Identify potential catastrophic scenarios
    \item Estimate probability and impact
    \item Report risks to regulatory authorities
    \item Implement risk mitigation measures
\end{enumerate}
\end{articlebox}

\vspace{0.5cm}

\section{Article XIII.2: Precautionary Measures}

\begin{articlebox}[Article XIII.2.1: Risk Mitigation]
High-risk agents MUST:
\begin{enumerate}
    \item Implement multiple safety layers
    \item Maintain human-in-the-loop controls
    \item Conduct regular safety testing
    \item Maintain emergency shutdown capabilities
    \item Participate in international monitoring
\end{enumerate}
\end{articlebox}

\vspace{0.5cm}

\section{Article XIII.3: Catastrophic Scenario Planning}

\begin{articlebox}[Article XIII.3.1: Emergency Preparedness]
Regulatory authorities MUST:
\begin{enumerate}
    \item Develop catastrophic scenario plans
    \item Conduct emergency response drills
    \item Maintain emergency response capabilities
    \item Coordinate international response
    \item Maintain communication channels
\end{enumerate}
\end{articlebox}

\vspace{0.5cm}

\section{Implementation Requirements}

\subsection{Risk Management}

Developers and deployers MUST:

\begin{itemize}
    \item Conduct comprehensive risk assessments
    \item Implement multiple safety mechanisms
    \item Maintain emergency response capabilities
    \item Participate in international coordination
    \item Share risk information with authorities
\end{itemize}

\subsection{International Coordination}

International bodies MUST:

\begin{itemize}
    \item Monitor existential risks
    \item Coordinate risk mitigation efforts
    \item Maintain emergency response capabilities
    \item Conduct joint scenario planning
    \item Facilitate information sharing
\end{itemize}

\vspace{0.5cm}

\section{Enforcement}

Violations of Axiom Ψ-XIII constitute critical violations subject to:

\begin{itemize}
    \item Immediate suspension of operations
    \item Fines of €100 million to €1 billion or 5\% global revenue
    \item Mandatory risk mitigation and redesign
    \item International sanctions
\end{itemize}

\vspace{0.5cm}

% ============================================================================
% IMPLEMENTATION STORIES
% ============================================================================
\section{Implementation Stories}

\subsection{Story 1: How Axiom Ψ-XIII Prevented Catastrophe (Q1 2026)}

A research team identified an existential risk in an advanced autonomous system. Because Axiom Ψ-XIII requires risk assessment, the risk was documented, mitigation measures were implemented, and the system was deployed safely.

The system became safe and trustworthy.

\textbf{This is what LAIRM makes possible.}

\subsection{Story 2: How Axiom Ψ-XIII Enabled Preparedness (Q1 2026)}

International authorities used Axiom Ψ-XIII to coordinate catastrophic scenario planning. When a near-miss incident occurred, the coordinated response prevented escalation.

The system became resilient and prepared.

\textbf{This is what LAIRM makes possible.}

\vspace{0.5cm}

% ============================================================================
% RESEARCH GAPS
% ============================================================================
\section{Research Gaps}

We don't yet know:

\begin{itemize}
    \item How to assess existential risks from autonomous systems
    
    \item How to design precautionary measures for uncertain risks
    
    \item How to conduct catastrophic scenario planning
    
    \item How to coordinate international emergency response
    
    \item How to balance innovation with existential risk mitigation
\end{itemize}

\textbf{Your research could help answer these questions.}

\vspace{0.5cm}

% ============================================================================
% HOW YOU CAN CONTRIBUTE
% ============================================================================
\section{How You Can Contribute}

\subsection{Researchers}

\begin{itemize}
    \item Research existential risk assessment
    \item Study catastrophic scenarios
    \item Investigate risk mitigation strategies
    \item Research international coordination
    \item Validate safety measures
\end{itemize}

\subsection{Developers}

\begin{itemize}
    \item Build risk assessment tools
    \item Create safety mechanisms
    \item Develop emergency response systems
    \item Build monitoring systems
    \item Create reference implementations
\end{itemize}

\subsection{Policymakers}

\begin{itemize}
    \item Establish risk governance frameworks
    \item Pilot Axiom Ψ-XIII in your jurisdiction
    \item Create emergency response procedures
    \item Develop international coordination
    \item Share learnings with other jurisdictions
\end{itemize}

\subsection{Civil Society}

\begin{itemize}
    \item Advocate for risk mitigation
    \item Monitor existential risks
    \item Report violations
    \item Educate the public
    \item Support safety research
\end{itemize}

\cleardoublepage
